\documentclass[11pt]{article}
\usepackage[a4paper,margin=28mm]{geometry}
\usepackage[T1]{fontenc}
\usepackage[utf8]{inputenc}
\usepackage{lmodern,microtype}
\usepackage{amsmath,amssymb,amsthm,mathtools}
\usepackage{enumitem}
\usepackage[dvipsnames]{xcolor}
\usepackage[colorlinks=true,linkcolor=MidnightBlue,citecolor=ForestGreen,urlcolor=MidnightBlue]{hyperref}
\hypersetup{pdftitle={Variational observables of log-correlated fields: a research plan for critical L1 transport}}
\setlist[enumerate]{itemsep=4pt,topsep=5pt}
\setlist[itemize]{itemsep=3pt,topsep=5pt}
\setlength{\emergencystretch}{2em}
\newtheorem{proposition}{Proposition}[section]
\theoremstyle{definition}
\newtheorem{target}[proposition]{Research target}
\newtheorem{remark}[proposition]{Remark}
\newcommand{\R}{\mathbb R}
\newcommand{\Z}{\mathbb Z}
\newcommand{\E}{\mathbb E}
\newcommand{\T}{\mathbb T}
\newcommand{\Lip}{\operatorname{Lip}}
\newcommand{\diver}{\operatorname{div}}
\newcommand{\curl}{\operatorname{curl}}
\newcommand{\Var}{\operatorname{Var}}
\newcommand{\Cov}{\operatorname{Cov}}
\newcommand{\KR}{\mathrm{KR}}
\newcommand{\cF}{\mathcal F}
\newcommand{\cM}{\mathcal M}
\newcommand{\PL}{P_{\mathrm L}}
\newcommand{\PT}{P_{\mathrm T}}
\newcommand{\dd}{\,dx}

\title{Variational Observables of Log-Correlated Fields\\[4pt]
\large A Research Plan for Critical $L^1$ Transport}
\author{Research plan notes}
\date{14 September 2026}

\begin{document}
\maketitle

\begin{abstract}
The proposed program studies the least total flux needed to realize the
divergence of a regularized log-correlated vector field. Its basic observable is
$\inf_{\diver K=0}\int|J_\varepsilon+K|$, normalized by
$\sqrt{\log(1/\varepsilon)}$. A Gaussian limit theorem for this observable,
together with quantitative approximation in the transport norm and control
of the boundary, would yield the limiting constant in Talagrand's
two-dimensional empirical $W_1$ problem. The relevant universality class is
determined by the longitudinal covariance of the field and by the transport
cost; logarithmic pointwise variance alone does not specify it. We formulate
concrete targets, give the exact normalization linking them to empirical
transport, and explain examples involving vector Gaussian free fields,
membrane slopes, random charge distributions, and network flows. We also
identify which connections to screening, stochastic homogenization, and
critical empirical processes are substantive, and which require additional
theorems. The statements labeled research targets are proposals, while the
elementary identities and conditional reductions are justified below.
\end{abstract}

\section{The question and the proposed generalization}
\label{sec:aim}

These notes develop the variational generalization proposed in the
discussion ``Prove Limit Existence.'' The hypotheses, conditional
reductions, and field examples below make that proposal precise.

Let $Q=[0,1]^2$, let $\lambda$ be Lebesgue probability measure on $Q$, and
let $X_1,X_2,\ldots$ be independent with law $\lambda$. Write
\[
  \mu_N=\frac1N\sum_{i=1}^N\delta_{X_i},\qquad
  A_N=\E\sup_{\Lip(f)\leq1}
  \left|\sum_{i=1}^N\left(f(X_i)-\int_Q f\dd\right)\right|.
\]
The critical order is $A_N\asymp\sqrt{N\log N}$. The question is whether
\begin{equation}\label{eq:talagrand}
  \frac{A_N}{\sqrt{N\log N}}
  =\sqrt{\frac N{\log N}}\,\E W_1(\mu_N,\lambda)
  \longrightarrow c_*\in(0,\infty).
\end{equation}
Talagrand explicitly formulates the limit problem in
\cite[Research Problem 4.5.4]{Talagrand}. The cited quadratic matching theory
\cite{AST} and the convergence results in dimensions $d\geq3$ \cite{GT}
give useful precedents. They do not provide the two-dimensional $p=1$
conclusion needed here. The nonlinear PDE reduction of \cite{AVT} concerns
strictly convex power costs $p>1$ and also distinguishes reduction to a
random PDE from existence of the asymptotic constant. These notes retain
\eqref{eq:talagrand} as a research target; the cited literature is the basis
for the program, rather than a claim of an exhaustive status survey.

The broader question suggested by the discussion is the following. Given
a smooth cutoff $J_\varepsilon$ of a random vector distribution on a domain
$D$, study
\begin{equation}\label{eq:informal}
  \cF_D(J_\varepsilon)
  :=\inf_{\diver K=0}\int_D|J_\varepsilon+K|\dd,
  \qquad
  \frac{\cF_D(J_\varepsilon)}{\sqrt{\log(1/\varepsilon)}}.
\end{equation}
The constraint on $K$ preserves sources and sinks while allowing the flux
to be rerouted. A boundary condition and a relaxed interpretation of the
integral are part of the definition, as specified next. The hoped-for
output is a deterministic limiting cost, together with an effective
description of the rerouting. This is a variational observable of a random
field, distinct from its maximum, its level sets, or an exponential measure.

\section{The exact transport dictionary}
\label{sec:dictionary}

\subsection{Balanced sources and fluxes}
For a finite signed measure $\sigma$ on a compact convex domain $D$, with
$\sigma(D)=0$, set
\[
  \|\sigma\|_{\KR(D)}
  =\sup_{\Lip(f)\leq1}\langle\sigma,f\rangle.
\]
The absolute value can be included without changing this supremum. When
$\sigma=\mu-\nu$ for probability measures, this equals $W_1(\mu,\nu)$.
The Beckmann formulation is
\begin{equation}\label{eq:beckmann}
  \|\sigma\|_{\KR(D)}
  =\inf\bigl\{|J|(\overline D):
       J\in\cM(\overline D;\R^2),\ -\diver J=\sigma,
       \ J\cdot n=0\bigr\}.
\end{equation}
Here the divergence and no-flux condition are imposed together by
\begin{equation}\label{eq:weakboundary}
  \int_{\overline D}\nabla f\cdot dJ=\int_{\overline D}f\,d\sigma
  \quad\text{for every }f\in C^1(\overline D).
\end{equation}
Thus no unmentioned boundary reservoir is allowed. Formula
\eqref{eq:beckmann} follows from transport duality and the minimal-flow
formulation; see \cite[Sections 3.1 and 4.2]{Santambrogio}. Convexity ensures
that the distance inside $D$ agrees with Euclidean distance. On a flat
torus we instead use periodic test functions and the torus distance.

If $J=J^a dx+J^s$, its cost is
$\int_D|J^a|\dd+|J^s|(\overline D)$. This relaxation accommodates atomic
sources and possible concentration of minimizing fluxes. We use the
same relaxation in \eqref{eq:informal}: $J_\varepsilon$ denotes a density,
and the actual cost is $|J_\varepsilon dx+K|(\overline D)$. For a norm
$\Phi$ on $\R^2$, write more generally
\[
  \|J\|_\Phi
  =\int_{\overline D}\Phi\!\left(\frac{dJ}{d|J|}\right)d|J|,
  \qquad
  \cF_D^\Phi(J_0)=\inf_{\diver K=0}\|J_0dx+K\|_\Phi.
\]
On a domain, $J_0$ and admissible $K$ also have zero normal flux. If
$\Phi^\circ$ is the dual norm, the dual constraint is
$\Phi^\circ(\nabla f)\leq1$. We write $\cF_D$ for $\Phi(v)=|v|$.

\subsection{Empirical fluctuations and the sign convention}
Define the balanced empirical fluctuation
\[
  \xi_N=\sqrt N(\mu_N-\lambda).
\]
Then, exactly,
\begin{equation}\label{eq:empiricalexact}
  \|\xi_N\|_{\KR(Q)}=\sqrt N\,W_1(\mu_N,\lambda),
  \qquad
  \frac{A_N}{\sqrt{N\log N}}
  =\frac{\E\|\xi_N\|_{\KR(Q)}}{\sqrt{\log N}}.
\end{equation}
Throughout, the sign convention is $-\diver J=\sigma$. Consequently,
if $-\Delta u=\sigma$ with Neumann boundary condition, the reference flux
is $J_0=\nabla u$, and every admissible flux is $J_0+K$ with
$\diver K=0$ and $K\cdot n=0$.

For fixed smooth test functions, the Gaussian limit of $\xi_N$ has
covariance
\begin{equation}\label{eq:bridge}
  \E\langle\xi,f\rangle\langle\xi,g\rangle
  =\int_Qfg\dd-\int_Qf\dd\int_Qg\dd.
\end{equation}
It is white noise with its constant mode removed: the spatial analogue
of a Brownian bridge. This finite-dimensional central limit statement
does not by itself control the growing supremum in
\eqref{eq:empiricalexact}. Section~\ref{sec:transfer} states the stronger
comparison that would suffice.

\section{Why the reference flux is log-correlated}
\label{sec:gaussian}

First work on the unit torus $\T^2=\R^2/\Z^2$ with volume one. Use Fourier
frequencies $k\in2\pi\Z^2$ and omit $k=0$. Let $\xi$ be mean-zero white
noise, let $\chi(r)=e^{-r^2/2}$, and define
\begin{equation}\label{eq:reference}
  \xi_\varepsilon=\chi(\varepsilon\sqrt{-\Delta})\xi,
  \qquad J_\varepsilon=\nabla(-\Delta)^{-1}\xi_\varepsilon.
\end{equation}
Its Fourier covariance is
\begin{equation}\label{eq:fourier}
  \E\bigl[\widehat J_\varepsilon(k)
      \widehat J_\varepsilon(k)^*\bigr]
  =\chi(\varepsilon|k|)^2\frac{k\otimes k}{|k|^4}.
\end{equation}
In particular, with $t_\varepsilon=\log(1/\varepsilon)$,
\begin{align}\label{eq:variance}
  \Cov(J_\varepsilon(x))
  &=\sum_{k\ne0}\chi(\varepsilon|k|)^2
        \frac{k\otimes k}{|k|^4}
    =\frac{t_\varepsilon}{4\pi}I_2+O(1).
\end{align}
The coefficient comes from the annular integral
\[
  \frac1{(2\pi)^2}\int_1^{1/\varepsilon}\frac{dr}{r}
            \int_{S^1}\theta\otimes\theta\,d\theta
  =\frac{\log(1/\varepsilon)}{4\pi}I_2.
\]
Replacing the lattice sum by this integral has bounded error: away from
the origin the derivative of the homogeneous summand is $O(|k|^{-3})$,
which is summable over unit cells in two dimensions; the smooth cutoff
adds a bounded error. The same Fourier kernel gives, at short separations,
\[
  \E[J_{\varepsilon,i}(x)J_{\varepsilon,j}(y)]
  =\frac{\delta_{ij}}{4\pi}
      \log\frac1{|x-y|\vee\varepsilon}+O(1),
\]
where the bounded remainder includes angular dependence. Thus
$O(t_\varepsilon)$ frequency annuli each contribute order-one variance.
This is the logarithmic scaling mechanism discussed in the general
Gaussian-field literature \cite{DRSV,LSSW}.

There is an explicit benchmark for the cost of the reference flux.
Stationarity, Gaussianity, and \eqref{eq:variance} imply
\begin{equation}\label{eq:linearbenchmark}
  \frac{\E\int_{\T^2}|J_\varepsilon|\dd}{\sqrt{t_\varepsilon}}
  \longrightarrow
  \E|Z|=\frac1{2\sqrt2},
  \qquad Z\sim N\!\left(0,\frac1{4\pi}I_2\right).
\end{equation}
Indeed, the expected radius of a two-dimensional standard Gaussian is
$\sqrt{\pi/2}$. The normalized Gaussian second moments are bounded, so
convergence in law also gives convergence of these first moments.
Taking $K=0$ yields the rigorous upper benchmark
\begin{equation}\label{eq:upper}
  \limsup_{\varepsilon\downarrow0}
  \frac{\E\cF_{\T^2}(J_\varepsilon)}{\sqrt{t_\varepsilon}}
  \leq\frac1{2\sqrt2}.
\end{equation}

The optimization remains essential. For $L^2$ fluxes, orthogonality gives
\[
  \int|\nabla u+K|^2
  =\int|\nabla u|^2+\int|K|^2
  \quad(\diver K=0),
\]
with periodic or compatible no-flux boundary conditions. There is no
such identity for $L^1$. Where $J_0\ne0$, a necessary stationarity
condition for a smooth $L^1$ minimizer is
\[
  \int\frac{J_0}{|J_0|}\cdot K\dd=0
  \quad\text{for all admissible solenoidal }K.
\]
Although $J_0=\nabla u$ is a gradient, $\nabla u/|\nabla u|$ generally is
not. Hence the Poisson flux need not minimize the cost. Whether the
reduction below \eqref{eq:linearbenchmark} survives at leading order is
one of the main structural questions.

\section{A precise class of general log-correlated vector fields}
\label{sec:general}

\subsection{The longitudinal covariance is the relevant datum}
Let $H_\varepsilon$ be a real, centered, stationary Gaussian vector field
on $\T^2$, with independent Fourier modes subject to the reality
constraint, and covariance
\begin{equation}\label{eq:generalcov}
  \E[\widehat H_\varepsilon(k)\widehat H_\varepsilon(k)^*]
  =\chi(\varepsilon|k|)^2\frac{A(k/|k|)}{|k|^2},
  \qquad k\ne0.
\end{equation}
Initially assume that $A$ is a smooth, even, real symmetric,
nonnegative matrix-valued function on $S^1$. Define
\[
  \PL(k)=\frac{k\otimes k}{|k|^2},\qquad
  \PT(k)=I_2-\PL(k),\qquad
  a(\theta)=\theta^TA(\theta)\theta.
\]
Every nonzero Fourier mode of $\PT H_\varepsilon$ is divergence-free.
Constant vector fields are divergence-free as well and are included in
the admissible corrections. Therefore
\begin{equation}\label{eq:quotient}
  \cF_{\T^2}^\Phi(H_\varepsilon)
  =\cF_{\T^2}^\Phi(\PL H_\varepsilon)
  =\| -\diver H_\varepsilon\|_{\KR,\Phi},
\end{equation}
where the last norm has dual constraint $\Phi^\circ(\nabla f)\leq1$.
The source $-\diver H_\varepsilon$ has Fourier variance
$\chi(\varepsilon|k|)^2a(k/|k|)$. Thus the full angular function $a$, rather
than the entire matrix $A$, determines the law of this observable.

\begin{remark}[Why logarithmic covariance alone is insufficient]
Let $R(v_1,v_2)=(-v_2,v_1)$. The field $RJ_\varepsilon$ is divergence-free,
so $\cF_{\T^2}(RJ_\varepsilon)=0$ identically. Nevertheless, it has the
same leading pointwise covariance $I_2t_\varepsilon/(4\pi)$ as
$J_\varepsilon$, and the same isotropic leading logarithmic singularity
in its covariance kernel. The angular part of the spectral covariance
distinguishes them. A positive universal limit for all fields having
only that logarithmic singularity would therefore be false.
\end{remark}

\begin{target}[Gaussian variational limit]\label{target:gaussian}
For every smooth even $a:S^1\to(0,\infty)$ with
$0<a_-\leq a\leq a_+<\infty$, and every norm $\Phi$ on $\R^2$, prove that
\begin{equation}\label{eq:gaussiantarget}
  \frac{\E\cF_{\T^2}^\Phi(H_\varepsilon)}{\sqrt{t_\varepsilon}}
  \longrightarrow\kappa(a,\Phi)\in(0,\infty).
\end{equation}
Prove independence of the limit from a specified class of smooth radial
cutoffs that equal one at zero and decay rapidly at infinity.
\end{target}
The first case to attack is $a\equiv1$, $\Phi=|\cdot|$; denote its
constant, if it exists, by $\kappa$. The general formulation isolates
anisotropic noise and anisotropic transport costs without assuming that
they give the same constant. Homogeneity forces
\[
  \kappa(s^2a,\Phi)=s\kappa(a,\Phi),\qquad
  \kappa(a,b\Phi)=b\kappa(a,\Phi)\quad(s,b>0).
\]
Strict positivity is part of the target. The upper bound follows directly
from the reference flux, while a lower bound must respect the global
gradient constraint in the dual problem.

\subsection{Stability and the mode of convergence}
Two elementary estimates help delimit the genuinely difficult part.
For fields on the same domain with the same boundary convention,
\begin{equation}\label{eq:stability}
  |\cF_D^\Phi(H)-\cF_D^\Phi(\widetilde H)|
  \leq\cF_D^\Phi(H-\widetilde H)
  \leq C_\Phi\int_D|H-\widetilde H|\dd.
\end{equation}
Consequently, a coupling with expected $L^1$ error
$o(\sqrt{t_\varepsilon})$ transfers the leading constant. In particular,
standard smooth cutoffs of the same Gaussian Fourier series have
bounded expected $L^1$ difference: their difference has variance
controlled by
\[
  C\int_1^\infty
       |\chi_1(\varepsilon r)-\chi_2(\varepsilon r)|^2\frac{dr}{r}
  \leq C'.
\]
This proves cutoff stability once existence is known. A bounded
pointwise difference of covariance kernels alone is not the coupling
estimate in \eqref{eq:stability}.

For the Gaussian source with bounded spectral density $a$, the variance
of the optimized cost is uniformly bounded:
\begin{equation}\label{eq:concentration}
  \Var\bigl(\cF_{\T^2}^\Phi(H_\varepsilon)\bigr)\leq C(a_+,\Phi).
\end{equation}
To see this, represent the source as a bounded Fourier multiplier of
white noise. Its dual cost is a supremum of Gaussian linear functionals
indexed by mean-zero $f$ with $\Phi^\circ(\nabla f)\leq1$.
Poincare's inequality bounds $\|f\|_2$ uniformly. This makes the
supremum uniformly Lipschitz in the underlying Gaussian coordinates;
the Gaussian Poincare inequality gives \eqref{eq:concentration}, first
for finitely many modes and then by approximation. Thus convergence of
the expectations in \eqref{eq:gaussiantarget} would automatically give
convergence in $L^2$ and in probability after normalization. It does not
follow from concentration that the normalized expectations converge.

\section{How a Gaussian theorem would imply Talagrand's limit}
\label{sec:transfer}

Return to $Q$ and use the Neumann heat semigroup $P_t$. Let
$S_\varepsilon=P_{\varepsilon^2/2}$, let $\xi$ have covariance
\eqref{eq:bridge}, and put
\[
  J_\varepsilon^Q
  =\nabla(-\Delta_{\mathrm N})^{-1}S_\varepsilon\xi,
  \qquad
  B_\varepsilon^Q
  =\cF_Q(J_\varepsilon^Q)
  =\|S_\varepsilon\xi\|_{\KR(Q)}.
\]
The inverse is on the mean-zero subspace. The no-flux condition is built
into this construction.

\begin{proposition}[A conditional transfer criterion]\label{prop:transfer}
Suppose that
\begin{equation}\label{eq:squarelimit}
  \frac{\E B_\varepsilon^Q}{\sqrt{\log(1/\varepsilon)}}
  \longrightarrow\kappa_Q\in(0,\infty).
\end{equation}
Suppose there are $\varepsilon_N\downarrow0$ such that
\begin{equation}\label{eq:scales}
  \frac{\log(1/\varepsilon_N)}{\log N}\longrightarrow\frac12,
  \qquad
  \sqrt N\,\varepsilon_N=o(\sqrt{\log N}),
\end{equation}
and couplings of $\xi_N$ and a field $\xi$ of law \eqref{eq:bridge} with
\begin{equation}\label{eq:couplingtarget}
  \E\|S_{\varepsilon_N}\xi_N-S_{\varepsilon_N}\xi\|_{\KR(Q)}
  =o(\sqrt{\log N}).
\end{equation}
Then \eqref{eq:talagrand} holds with $c_*=\kappa_Q/\sqrt2$.
\end{proposition}

\begin{proof}
Reflected Brownian motion on the square has expected displacement at
time $\varepsilon^2/2$ bounded by $C\varepsilon$, uniformly in its
starting point. Coupling a point to its heat evolution gives
$W_1(\nu,S_\varepsilon\nu)\leq C\varepsilon$ for every probability
measure $\nu$. Since $S_\varepsilon\lambda=\lambda$,
\[
  \|\xi_N-S_\varepsilon\xi_N\|_{\KR(Q)}
  \leq C\sqrt N\,\varepsilon.
\]
The triangle inequality and \eqref{eq:couplingtarget} therefore imply
\[
  \left|\E\|\xi_N\|_{\KR(Q)}-\E B_{\varepsilon_N}^Q\right|
  =o(\sqrt{\log N}).
\]
Now use \eqref{eq:squarelimit}, \eqref{eq:scales}, and
\eqref{eq:empiricalexact}.
\end{proof}

A useful scale to test is
\begin{equation}\label{eq:mesoscale}
  \varepsilon_N=N^{-1/2}(\log N)^\beta,
  \qquad 0<\beta<\frac12.
\end{equation}
The expected number of sample points in a cell of that size is
$N\varepsilon_N^2=(\log N)^{2\beta}\to\infty$, while the smoothing error
is negligible at the target scale. This leaves room for a cellwise
Gaussian approximation without discarding a positive fraction of the
logarithmic scales. The coupling in \eqref{eq:couplingtarget} remains a
quantitative task; a central limit theorem for each fixed test function
does not prove it.

One possible route is to Poissonize the cell counts, couple each
centered count to a Gaussian, balance the resulting residual, and bound
its transport cost through an inverse-divergence estimate. For a smooth
mean-zero residual $r$,
\begin{equation}\label{eq:hminusone}
  \|r\|_{\KR(Q)}
  \leq\|\nabla(-\Delta_{\mathrm N})^{-1}r\|_{L^1(Q)}
  \leq\|r\|_{H^{-1}_{\mathrm N}(Q)}.
\end{equation}
The last norm is defined by the displayed Neumann gradient in $L^2$.
Cellwise residuals must be spread over cells or smoothed before this
estimate is used: point charges in dimension two do not have finite
$H^{-1}$ norm. The errors from cell placement, Gaussian coupling,
balancing, and removal of Poissonization must all satisfy the same
$o(\sqrt{\log N})$ budget.

There is a separate geometric step. A theorem on $\T^2$ does not
automatically imply \eqref{eq:squarelimit}; periodic and no-flux
admissible classes differ. A boundary comparison must prove
$\kappa_Q=\kappa$. With that comparison and
Proposition~\ref{prop:transfer}, the constants are related by
\begin{equation}\label{eq:constants}
  c_*=\frac\kappa{\sqrt2},\qquad
  \kappa\leq\frac1{2\sqrt2},\qquad c_*\leq\frac14.
\end{equation}
The last inequality is the consequence of these transfer inputs and the
explicit Gaussian benchmark; equality is a further question about the
optimal solenoidal correction.

\section{Concrete fields and their mathematical or physical meaning}
\label{sec:examples}

\subsection{Two independent GFFs: an exact Gaussian identification}
Let $h_1,h_2$ be independent mean-zero Gaussian free fields on $\T^2$,
each with covariance operator $(-\Delta)^{-1}$, and set
$H_\varepsilon=(S_\varepsilon h_1,S_\varepsilon h_2)$. Then
\[
  \E[\widehat H_\varepsilon(k)\widehat H_\varepsilon(k)^*]
  =\chi(\varepsilon|k|)^2\frac{I_2}{|k|^2}.
\]
Its longitudinal projection has covariance
$\chi(\varepsilon|k|)^2(k\otimes k)/|k|^4$. Hence, as Gaussian fields,
\begin{equation}\label{eq:vectorgff}
  \PL H_\varepsilon\ \stackrel{\mathrm{law}}=\ J_\varepsilon,
  \qquad
  \cF_{\T^2}(H_\varepsilon)
     \ \stackrel{\mathrm{law}}=\ \cF_{\T^2}(J_\varepsilon).
\end{equation}
This supplies a concrete observable of a familiar free field with
exactly the same Gaussian variational problem. In the free-field
interpretation, the two components have quadratic gradient energy
$\frac12\int(|\nabla h_1|^2+|\nabla h_2|^2)$, understood with a cutoff.
The proposed observable measures the least total vector amplitude
compatible with the field's sources and sinks. Its value forgets the
transverse component, even though that component contributes to the
ordinary $L^1$ norm. This is a precise connection to GFF theory, but it
does not yet produce a statement about GFF extremes or multiplicative
chaos.

\subsection{Membrane slopes: the reference field itself}
Let
\[
  u=(-\Delta)^{-1}\xi,
  \qquad \Cov(u)=(-\Delta)^{-2}.
\]
This is the mean-zero bi-Laplacian Gaussian field. Its discrete
counterpart is the Gaussian membrane model \cite{LSSW}. At a spectral
cutoff it has the quadratic bending Hamiltonian
\begin{equation}\label{eq:membrane}
  \mathcal H(u)=\frac12\int_{\T^2}(\Delta u)^2\dd,
  \qquad J_\varepsilon=\nabla u_\varepsilon.
\end{equation}
In the small-slope approximation to a tensionless membrane, $u$ is a
height and $\nabla u$ is its slope. In dimension two, it is this slope
field that is log-correlated. The height is not the ordinary
two-dimensional GFF. In dimension four the bi-Laplacian height itself
is log-correlated, which is a different instance of the fractional
Gaussian-field family.

The functional $\cF(\nabla u_\varepsilon)$ asks how much of the integrated
slope magnitude remains after arbitrary divergence-free corrections
while retaining $-\Delta u_\varepsilon$. This is a mathematically natural
nonlinear observable of membrane slopes. Such corrections need not be
gradients of another height, so the optimization is not itself the
physical bending energy or a membrane equilibrium principle. A useful
result would quantify the difference between unconstrained slope
amplitude and the amplitude forced by the curvature source.

\subsection{Random charges: an \texorpdfstring{$L^1$}{L1} counterpart of electrostatics}
For a smooth neutral charge density $\rho_\varepsilon$, let
$-\Delta V_\varepsilon=\rho_\varepsilon$. The electrostatic field
$E_\varepsilon=-\nabla V_\varepsilon$ satisfies
$\diver E_\varepsilon=\rho_\varepsilon$. Our reference flux is
$J_\varepsilon=-E_\varepsilon$. In the quadratic variational principle,
the gradient field minimizes $\frac12\int|E|^2$ among fields with that
divergence. The corresponding linear-growth problem is
\begin{equation}\label{eq:charge}
  \inf_{\diver E=\rho_\varepsilon}\int|E|\dd
  =\|\rho_\varepsilon\|_{\KR}.
\end{equation}
For an ordinary signed density this is exactly the transport cost from
its positive part to its negative part, which have equal total mass.
It measures the work needed to neutralize the charge with unit cost per
unit mass and distance.

This distinction gives a concrete purpose to screening estimates:
neutralize charge locally, match the remaining boundary fluxes, and
bound the extra transport required to join neighboring regions. The
charge is prescribed in \eqref{eq:charge}; no additional Coulomb-gas
Gibbs weight is implicit.

The spectral density also explains why Coulomb systems should not all
be placed in the white-noise universality class. In a large box of side
$L$, with a fixed microscopic cutoff, a stationary isotropic charge
field of structure factor $S(k)$ has a linear-field variance contribution
proportional to
\begin{equation}\label{eq:structurefactor}
  \int_{1/L}^{1}\frac{S(r)}r\,dr.
\end{equation}
If $S(r)\to S_0>0$, this is $S_0\log L+o(\log L)$. If instead
$S(r)=O(r^\alpha)$ with $\alpha>0$, this contribution stays bounded.
The latter behavior is associated with hyperuniform suppression of
long-wavelength density fluctuations \cite{TS}. Formula
\eqref{eq:structurefactor} is a diagnostic for the linear flux, not a
complete asymptotic theorem for the optimized cost. It identifies
screened charge systems as tests of where the critical normalization
changes.

\subsection{Random network flows and non-Gaussian noise}
On a finite connected graph with oriented incidence matrix $B$, neutral
charges $\eta=(\eta_v)$, and positive edge lengths $\ell_e$, the exact
finite-dimensional analogue is
\begin{equation}\label{eq:graph}
  \min_{Bj=\eta}\sum_e\ell_e|j_e|
  =\max_{| (B^Tf)_e|\leq\ell_e}\sum_v\eta_vf_v.
\end{equation}
This is linear programming duality. The observable is minimum routing
cost for random supply and demand. Independent centered Bernoulli,
Poisson, and Gaussian charges provide distinct microscopic laws to
compare. The mean charge must be removed, or otherwise balanced,
before solving the problem.

In a continuum box of side $L$ at unit source intensity, the Gaussian
critical prediction is total cost of order $L^2\sqrt{\log L}$.
For the empirical model with $N=L^2$, scaling space by $L$ and mass by
$N$ gives the exact relation
\[
  W_1\!\left(\sum_{i=1}^N\delta_{LX_i},
                   \mathcal L^2\!\restriction_{[0,L]^2}\right)
  =L^3W_1(\mu_N,\lambda).
\]
Here $W_1$ is extended to equal finite masses. Thus
\eqref{eq:talagrand} would give coefficient $\sqrt2c_*=\kappa$ in front
of $L^2\sqrt{\log L}$, provided the boundary identification holds.

A nearest-neighbor square graph uses the path metric and has an
anisotropic continuum cost. Its limiting coefficient should therefore
be compared with $\kappa(a,\Phi)$ for the appropriate $\Phi$, not directly
with the Euclidean $\kappa$. A numerical Euclidean discretization must
approximate the Euclidean norm and the divergence operator consistently.

For stationary lattice charges of finite dependence range, a natural
candidate noise strength is
\[
  s^2=\sum_{z\in\Z^2}\Cov(\eta_0,\eta_z).
\]
With adequate moments, $s^2>0$, and a quantitative transport comparison,
one may seek coefficient $s\kappa$ for an isotropic continuum cost.
Finite variance and short-range dependence alone are not asserted here
to prove this universality. Even finite-range noise can have $s^2=0$:
for example, $\eta_z=\zeta_z-\zeta_{z-e_1}$ for independent centered
$\zeta_z$. On a periodic graph this example can be transported by a local edge flux and
does not have the same positive critical coefficient.

\subsection{Two further tests of the formulation}
\paragraph{A scalar GFF used as the source.}
Let $h$ be a mean-zero scalar GFF on $\T^2$. The transport norm of
$h_\varepsilon$ is a different problem from
\eqref{eq:vectorgff}. Indeed,
\[
  \E\|\nabla(-\Delta)^{-1}h\|_2^2
  =\sum_{k\ne0}|k|^{-4}<\infty.
\]
Hence $\|h_\varepsilon\|_{\KR}$ has bounded expectation, and in fact
converges in $L^1$ as the cutoff is removed, by the $L^2$ convergence of
its reference flux and \eqref{eq:stability}. Dividing by
$\sqrt{\log(1/\varepsilon)}$ gives zero. One must specify whether a
log-correlated field is the source, the flux, or a potential.

\paragraph{A curl constraint.}
For a vector field $H$, define
$\cF_{\curl}(H)=\inf_{\curl K=0}\int|H+K|$ with the same relaxed
interpretation on the torus. Since $\diver(R^TH)=\curl H$ and rotations
preserve Euclidean length,
\[
  \cF_{\curl}(H)=\cF_{\T^2}(R^TH).
\]
Consequently, $RJ_\varepsilon$, which is trivial for the divergence
observable, gives the original problem for the curl observable. This
describes minimum integrated velocity magnitude at prescribed vorticity
when arbitrary curl-free corrections are allowed. Imposing
incompressibility as an additional constraint would produce a different
problem. This example suggests studying a field together with its
differential constraint, rather than assigning a variational theory to
the covariance alone.

\section{Domains, densities, and extensions of the field class}
\label{sec:extensions}

For a probability density $\rho$ on a bounded convex planar domain $D$,
assume initially $\rho$ is smooth and bounded above and below by positive
constants. The empirical Gaussian source can be realized as
\begin{equation}\label{eq:densitybridge}
  \xi_\rho=\sqrt\rho\,W
       -\rho\int_D\sqrt\rho\,dW,
\end{equation}
where $W$ is white noise on $D$. Its covariance is
$\int_D\rho fg-\int_D\rho f\int_D\rho g$ and its total mass is zero.
Freezing the density near $x$ gives local noise amplitude
$\sqrt{\rho(x)}$. The rank-one balancing term is smooth under these
assumptions and has bounded expected transport cost. The prospective
locality formula is therefore
\begin{equation}\label{eq:densityprediction}
  \sqrt{\frac N{\log N}}\,\E W_1(\mu_N,\rho dx)
  \longrightarrow c_*\int_D\sqrt{\rho(x)}\dd.
\end{equation}
This is a prediction conditional on localization, Gaussian comparison,
and boundary control. The power $\sqrt\rho$ follows from noise amplitude
and the one-homogeneity of the cost; it is not imported from the
quadratic transport formula. Under dilation of $D$ by a factor $r$,
the right side scales by $r$, as a transport distance should. Densities
that vanish or are unbounded require separate analysis.

More generally, a slowly varying longitudinal symbol $a_x(\theta)$ and
cost $\Phi_x$ suggest a Gaussian bulk expression
\[
  \int_D\kappa(a_x,\Phi_x)\dd.
\]
This is an additional local limit target, beyond the stationary
problem. Changing the elliptic operator used to generate a field does
not by itself identify the transport constant: one must calculate its
resulting longitudinal symbol.

The spectral idea also makes sense in dimension $d$: a vector field with
covariance $|k|^{-d}A(k/|k|)$ is logarithmically critical, and its
divergence has source spectrum $|k|^{2-d}a(k/|k|)$. White noise is this
source precisely when $d=2$. Thus a higher-dimensional log-correlated
field extension would concern colored sources and should not be
confused with the ordinary high-dimensional empirical matching problem.

\section{A proof program organized around the missing estimates}
\label{sec:program}

\subsection{Preserve the coarse flux when combining scales}
The primary difficulty is that optimizing each annulus separately does
not optimize the sum. If $J_\varepsilon=\sum_{j\leq m}J^{(j)}$, triangle
inequalities can charge each of $m\asymp t_\varepsilon$ scales
separately, giving an order-$m$ bound and losing the order-$\sqrt m$
behavior. Even when a sharper upper bound is available, scalar costs
alone do not retain the directional information needed for a limiting
constant.

A plausible state variable is a convex effective cost depending on the
imposed mean flux. For the mean-zero reference field on $\T^2$, define
\begin{equation}\label{eq:profile}
  F_\varepsilon(q,\omega)
  =\frac1{\sqrt{t_\varepsilon}}
      \inf_{\substack{\diver K=0\\\int_{\T^2}dK=0}}
       \left\|\sqrt{t_\varepsilon}\,q\,dx
                      +J_\varepsilon dx+K\right\|_\Phi,
  \qquad q\in\R^2.
\end{equation}
The mean-zero condition on $K$ is indispensable here: without it the
constant vector $q$ could be canceled. The unrestricted observable is
exactly
\[
  \frac{\cF_{\T^2}^\Phi(J_\varepsilon)}{\sqrt{t_\varepsilon}}
  =\inf_{q\in\R^2}F_\varepsilon(q,\omega).
\]
The profiles are convex and satisfy
$F_\varepsilon(q,\omega)\geq\Phi(q)$ by Jensen's inequality for vector
measures. Their expectations are even in $q$ by Gaussian symmetry.
If one proves locally uniform concentration and convergence to a
deterministic profile $\overline\Phi_a$, together with control of
minimizing $q$, then symmetry and convexity identify
$\kappa(a,\Phi)=\overline\Phi_a(0)$. Interchanging expectation and the
infimum before this control would be an extra, unjustified step.

The intended multiscale operation is to absorb a band of finer scales
into such a convex cost, conditional on the coarse flux. A successful
construction must justify spatial localization, boundary matching, and
the sum of the errors over all bands. Any limiting flow equation for
the profiles should be derived from that operation. Logarithmic
variance or an informal renormalization analogy does not determine its
generator.

\subsection{Boundary screening and a lower-bound mechanism}
An upper construction can join local transports by correcting their
boundary imbalance. A matching lower bound must show that a global
optimizer cannot exploit the artificial interfaces to gain an amount
of order $\sqrt{t_\varepsilon}$. In dual language, local potentials
must be joined while preserving $\Phi^\circ(\nabla f)\leq1$ and losing
only $o(\sqrt{t_\varepsilon})$ in source pairing. In primal language,
one needs control of boundary traces and the cost of their correction.
The same estimates are the likely route from a torus to the square
and to \eqref{eq:densityprediction}.

This is the concrete sense in which the project may develop a critical
stochastic homogenization theory: it seeks a deterministic constitutive
cost after optimizing over logarithmically many active scales. It is
not an immediate application of a conventional fixed-scale cell
problem or a subadditive theorem. In particular, stabilization of
increments would have to be formulated for a suitable profile or
renormalized quantity; the raw $L^1$ costs are not additive in scales.

\subsection{The endpoint of power-cost PDEs}
For a fixed smooth source $\sigma$, consider the static flux problem
\[
  B_p(\sigma)=\inf_{-\diver J=\sigma}\frac1p\int|J|^p\dd,
  \qquad p>1.
\]
Its convex dual is
\[
  \sup_f\left\{\langle\sigma,f\rangle
                   -\frac1q\int|\nabla f|^q\dd\right\},
  \qquad q=\frac p{p-1},
\]
and formally the maximizer solves
$-\diver(|\nabla f|^{q-2}\nabla f)=\sigma$. As $p\downarrow1$,
$q\uparrow\infty$ and the dual penalty approaches the hard constraint
$|\nabla f|\leq1$. This supplies a useful interpretation of the
degeneracy at $p=1$. The static problem $B_p$ is not an exact formula
for $W_p^p$ when $p>1$; its link to matching requires a PDE
approximation such as the one developed in \cite{AVT}.

An endpoint route must control $p\downarrow1$ uniformly as the cutoff
vanishes. Convergence of the deterministic problem at each fixed
cutoff would not justify interchanging these two limits. Possible
concentration of the primal flux is another reason to keep the
measure-valued formulation.

\subsection{Priority order and concrete deliverables}
\begin{enumerate}[label=\textbf{\arabic*.}]
  \item \textbf{Gaussian base case.} Establish sharp two-sided order
  for the relaxed white-noise problem, then construct a mechanism
  proving convergence of its normalized expectation. Track the gain
  over the reference-flux benchmark \eqref{eq:linearbenchmark}.
  \item \textbf{Effective profile and boundary comparison.} Prove
  compactness, concentration, and identification for
  \eqref{eq:profile}, or replace it by another explicitly defined
  state variable that closes under scale elimination. Prove the
  torus--square comparison with a quantitative error budget.
  \item \textbf{Empirical transfer.} Establish
  \eqref{eq:couplingtarget} at a scale such as \eqref{eq:mesoscale}.
  Proposition~\ref{prop:transfer} then records exactly how the
  Gaussian result yields Talagrand's constant.
  \item \textbf{Universality and locality.} Extend first to independent
  non-Gaussian cell charges with specified moment assumptions, then
  to finite-range dependence with positive integrated covariance.
  Treat smooth densities and anisotropic symbols using the same
  localization estimates.
  \item \textbf{Field-specific consequences.} Translate the Gaussian
  theorem to \eqref{eq:vectorgff} and \eqref{eq:membrane}; these require
  no empirical approximation. Study hyperuniform sources and altered
  differential constraints as separate extensions, allowing their
  normalization to change.
\end{enumerate}

Numerical work can test the first and second deliverables before a full
proof is available. A consistent convex flux discretization should
record the optimized cost, the reference-flux cost, their difference,
and the primal--dual gap across several cutoffs. Vector GFFs and their
longitudinal projections provide an exact-law check; a purely
solenoidal input provides a zero-cost check. Comparisons across noise
laws should match integrated covariance and transport norm. Apparent
convergence over a short range of logarithmic scales is evidence for
a conjecture, not a substitute for an error estimate.

\section{What a successful result would explain}
\label{sec:meaning}

For random optimal transport, the immediate output is a first-order
constant and an account of whether solenoidal rerouting changes it.
For log-correlated fields, the Gaussian theorem would provide a
nonlinear observable with a deterministic bulk asymptotic, and would
show precisely which part of the covariance survives the optimization.
The vector-GFF and membrane identities make this consequence concrete.
For random charge and network models, it would describe minimum
neutralization or routing cost and test its dependence on microscopic
noise. The boundary estimates and effective profiles could then become
tools for other critical stochastic variational problems.

There is also an empirical-process direction. In the dual description,
one is seeking exact asymptotics for a Gaussian or empirical supremum
over a critical Lipschitz class. Extensions to H\"older, Sobolev, or
Besov test classes would require their own dual constraints and
normalizations. A matching entropy exponent alone cannot give the
constant, or even guarantee the same logarithmic correction: a direct
entropy integral can lose a logarithmic factor.

Finally, a theorem about this spatially integrated convex observable
would not by itself identify a maximum, an extremal process, a
derivative martingale, or a Gaussian multiplicative chaos measure.
Connections to those observables require an additional argument
involving high points or multiplicative structure. The attainable
central objective is more specific:
\[
  \boxed{\begin{gathered}
    \text{A variational limit for the longitudinal part of a}\\
    \text{log-correlated field, with quantitative transport universality.}
  \end{gathered}}
\]
Within that program, Talagrand's empirical problem is the application
that additionally requires a Gaussian approximation and a boundary
comparison at the scale of the leading constant.

\begin{thebibliography}{99}

\bibitem{Talagrand}
M.~Talagrand,
\emph{Upper and Lower Bounds for Stochastic Processes: Decomposition
Theorems}, second edition, Springer, 2021.
Section~4.5, especially Research Problem~4.5.4.
\href{https://michel.talagrand.net/ULBSPRINGER.pdf}{Author's PDF};
\href{https://doi.org/10.1007/978-3-030-82595-9}{published edition}.

\bibitem{Santambrogio}
F.~Santambrogio,
\emph{Optimal Transport for Applied Mathematicians: Calculus of
Variations, PDEs, and Modeling}, Birkh\"auser, 2015.
Sections~3.1 and 4.2.
\href{https://doi.org/10.1007/978-3-319-20828-2}{DOI: 10.1007/978-3-319-20828-2}.

\bibitem{AST}
L.~Ambrosio, F.~Stra, and D.~Trevisan,
\emph{A PDE approach to a 2-dimensional matching problem},
Probability Theory and Related Fields \textbf{173} (2019), 433--477.
\href{https://arxiv.org/abs/1611.04960}{arXiv:1611.04960};
\href{https://doi.org/10.1007/s00440-018-0837-x}{published article}.

\bibitem{GT}
M.~Goldman and D.~Trevisan,
\emph{Convergence of asymptotic costs for random Euclidean matching
problems}, 2020 preprint.
\href{https://arxiv.org/abs/2009.04128}{arXiv:2009.04128}.

\bibitem{AVT}
L.~Ambrosio, F.~Vitillaro, and D.~Trevisan,
\emph{Sharp PDE estimates for random two-dimensional bipartite matching
with power cost function}, Rendiconti Lincei--Matematica e Applicazioni
\textbf{35} (2024), 323--342.
\href{https://arxiv.org/abs/2405.09397}{arXiv:2405.09397};
\href{https://doi.org/10.4171/RLM/1042}{published article}.

\bibitem{DRSV}
B.~Duplantier, R.~Rhodes, S.~Sheffield, and V.~Vargas,
\emph{Log-correlated Gaussian fields: an overview}, 2014 preprint.
\href{https://arxiv.org/abs/1407.5605}{arXiv:1407.5605}.

\bibitem{LSSW}
A.~Lodhia, S.~Sheffield, X.~Sun, and S.~S.~Watson,
\emph{Fractional Gaussian fields: a survey}, 2014 preprint.
\href{https://arxiv.org/abs/1407.5598}{arXiv:1407.5598}.

\bibitem{TS}
S.~Torquato and F.~H.~Stillinger,
\emph{Local density fluctuations, hyperuniformity, and order metrics},
Physical Review E \textbf{68} (2003), 041113.
\href{https://doi.org/10.1103/PhysRevE.68.041113}{DOI: 10.1103/PhysRevE.68.041113}.

\end{thebibliography}
\end{document}
